\section{One-electron operators and GTH pseudopotential derivatives}
\label{sec:gth-one-electron}

\subsection{The periodic one-electron Hamiltonian}

With a GTH pseudopotential, the one-electron Hamiltonian is
\begin{equation}
  \hat h
  =-\frac12\nabla^2
   +\sum_A\hat V_{\mathrm{loc}}^A
   +\sum_A\hat V_{\mathrm{nl}}^A.
  \label{eq:gth-hamiltonian}
\end{equation}
The all-electron electron--nuclear potential $-Z_A/r$ is replaced by a smooth
local potential and a finite-rank nonlocal projector.  Although the kinetic
operator has no explicit nuclear-coordinate dependence, its AO matrix
derivative is nonzero because the atom-centered basis moves.

\subsection{The local GTH potential}

A representative local form is
\begin{equation}
  V_{\mathrm{loc}}^A(r)
  =-\frac{Z_A^{\mathrm{ion}}}{r}
    \operatorname{erf}\!\left(\frac{r}{\sqrt2 r_{\mathrm{loc},A}}\right)
   +e^{-r^2/(2r_{\mathrm{loc},A}^2)}
    \sum_{n=1}^{4}C_{nA}
    \left(\frac{r}{r_{\mathrm{loc},A}}\right)^{2n-2}.
  \label{eq:gth-local-real}
\end{equation}
In a periodic FFT representation, the position dependence is carried by the
structure factor $e^{-i\vG\cdot\mathbf R_A}$, and therefore
\begin{equation}
  \frac{\partial}{\partial R_{Ax}}
  \left[e^{-i\vG\cdot\mathbf R_A}V_{\mathrm{loc}}^A(\vG)\right]
  =-iG_xe^{-i\vG\cdot\mathbf R_A}V_{\mathrm{loc}}^A(\vG).
  \label{eq:gth-local-structure-derivative}
\end{equation}
At $\vG=0$, $G_x=0$, but the constant term, charge convention, and Ewald
compensation must still match those used by the central hcore.

The AO matrix is
\begin{equation}
  V_{\mu\nu,\mathrm{loc}}^{\vk}
  =\langle\phi_{\mu\vk}|\hat V_{\mathrm{loc}}|\phi_{\nu\vk}\rangle.
  \label{eq:gth-local-ao}
\end{equation}
Its derivative contains the motion of potential center $A$, the bra AO center,
and the ket AO center.  Translational identities can express the
potential-center derivative through a combination of AO derivatives and the
reciprocal-space structure-factor derivative.

\subsection{Nonlocal projectors}

The nonlocal GTH operator has a separable form for each angular-momentum
channel,
\begin{equation}
  \hat V_{\mathrm{nl}}^A
  =\sum_{l m}\sum_{ij}
   |p_{i,lm}^{A}\rangle h_{ij}^{A,l}
   \langle p_{j,lm}^{A}|.
  \label{eq:gth-nonlocal-operator}
\end{equation}
Using the projector overlap
\begin{equation}
  P_{i,lm;\mu}^{A\vk}
  =\langle p_{i,lm}^{A}|\phi_{\mu\vk}\rangle,
  \label{eq:gth-projector-overlap}
\end{equation}
the AO matrix element is
\begin{equation}
  V_{\mu\nu,\mathrm{nl}}^{\vk}
  =\sum_{A,lm,ij}
   P_{i,lm;\mu}^{A\vk *}h_{ij}^{A,l}
   P_{j,lm;\nu}^{A\vk}.
  \label{eq:gth-nonlocal-matrix}
\end{equation}
The product rule gives
\begin{align}
  V_{\mu\nu,\mathrm{nl}}^{\vk,Ax}
  =\sum_{l m i j}
  \bigg[&P_{i,lm;\mu}^{A\vk,Ax *}h_{ij}^{A,l}
         P_{j,lm;\nu}^{A\vk}
  \notag\\
       +&P_{i,lm;\mu}^{A\vk *}h_{ij}^{A,l}
         P_{j,lm;\nu}^{A\vk,Ax}\bigg]
  +\text{AO-center terms}.
  \label{eq:gth-nonlocal-derivative}
\end{align}
The radial projector table $h_{ij}^{A,l}$ moves rigidly with the atom and has
no parameter derivative with respect to the nuclear coordinate.

\subsection{AO-row conventions and the full matrix}

An integral engine often returns a derivative matrix $h_{\mu\nu}^{(1)}$ with
respect to the AO-row center.  The full matrix derivative for atom $A$ must
combine bra, ket, and operator-center ownership,
\begin{equation}
  h^{Ax}=h_{\mathrm{row},A}^{x}
       +h_{\mathrm{row},A}^{x\dagger}
       +h_{\mathrm{operator},A}^{x}.
  \label{eq:hcore-full-derivative}
\end{equation}
At complex k-points, the second term is a conjugate transpose, not merely a
transpose.

The overlap derivative obeys the analogous AO-center identity
\begin{equation}
  S^{Ax}=S_{\mathrm{row},A}^{x}+S_{\mathrm{row},A}^{x\dagger}.
  \label{eq:overlap-full-derivative}
\end{equation}
Both $h^{Ax}$ and $S^{Ax}$ enter the explicit target gradient and the source
HF CPHF source in Eq.~\eqref{eq:cphf-source}.

\subsection{Pseudopotential consistency}

Using an HF source density does not create a distinct analytic form called an
``HF pseudopotential.''  A label such as GTH-HF identifies a pseudopotential
parameter set.  The source HF calculation and the target-DFA energy and
gradient must use the same GTH parameters and valence ionic charges.  Using
different pseudopotentials for source and target would no longer correspond to
differentiating a single energy surface of Eq.~\eqref{eq:scope-dc-energy}.

Validation first compares local, nonlocal, kinetic/overlap, and ion--ion
components separately with finite differences and then checks the full hcore
derivative.  Atom-slice and conjugation errors become especially visible when
the nonlocal projector contains more angular channels.
